\documentclass[protokoll]{dvd}

\KOMAoptions{
    headwidth = 18cm,
    footwidth = 18cm,
}

% --------------------------------------------------------------------- %
% ---------------------------- Definitions ---------------------------- %
% --------------------------------------------------------------------- %

% Styrelsen
\def\ordf{Namn Ordförande}
\def\vice{Namn Vice}
\def\samo{Namn SAMO}
\def\kassa{Namn Kassör}
\def\sekr{Namn Sekreterare}
\def\vak{\emph{vakant}}

% mötesinfo

\def\nr{X}
\def\yr{yyyy}
\def\ymd{YYYY-MM-DD}
\def\starttime{HH:MM}

% --------------------------------------------------------------------- %

\begin{document}

\title{Styrelsemöte \nr}
\subtitle{\yr}
\author{Styrelsen}
\date{\ymd}


\textbf{Datum:} \csname @date\endcsname\\
\textbf{Tid:} \starttime\\
\textbf{Plats:} Styrelserummet\\
\textbf{Styrelsemedlemmar:}
\begin{närvarande_förtroendevalda}
    \förtroendevald{Ordförande}{\ordf}{Ja/Nej}
    \förtroendevald{\vice}{Tim Persson}{Ja/Nej}
    \förtroendevald{\kassa}{Lukas Gartman}{Ja/Nej}
    \förtroendevald{\samo}{Josefin Kokkinakis}{Ja/Nej}
    \förtroendevald{\sekr}{Gustav Dalemo}{Ja/Nej}
\end{närvarande_förtroendevalda}

%\textbf{Övriga medlemmar:}
%\begin{närvarande_medlemmar}
%\medlem{Ingen}[]
%\end{närvarande_medlemmar}
\section{Öppnande av möte}

Mötet öppnades av Samuel Hammersberg kl 11:13

\section{Runda bordet}

Runda bordet innebär att varje person berättar hur de känner sig.
Man kan till exempel berätta att man är stressad på grund av en inlämning, irriterad på sin granne, eller bara väldigt glad därför att man ligger i fas med plugget.

\section{Formalia}

\subsection{Styrelsens beslutbarhet}

\blockquote[7 kap. 5 \S~första stycket i stadgan][]{%
    Styrelsen är endast beslutsmässig då samtliga styrelsemedlemmar har fått kallelsen till styrelsemötet och minst hälften av styrelsemedlemmarna är närvarande.
    Ordförande eller vice ordförande måste vara närvarande när beslut tas.
}

\subsubsection*{Förslag}

\begin{attsatser}
    \item Styrelsen har uppnått kraven i 7 kap. 5 § första stycket i stadgan och är därmed beslutbar.
\end{attsatser}
\subsubsection*{Beslut}
\begin{attsatser}
    \item %förslaget till beslut bifalles
\end{attsatser}


\subsection{Fastställande av mötesschema}

För att styrelsen ska kunna fatta ett styrelsebeslut eller protokollföra en diskussion behöver punkten i mötesschemat där styrelsen ska fatta beslut vara inlagd eller föras in i mötesschemat senast vid den här punkten.

\subsubsection*{Förslag}

\begin{attsatser}
    \item 
    %mötesschemat fastställs utan några förändringar.
    %punkt läggs till på mötesschema
\end{attsatser}
\subsubsection*{Beslut}
\begin{attsatser}
    \item %förslaget till beslut bifalles
\end{attsatser}

\subsection{Val av mötessekreterare}
...

\subsubsection*{Förslag}
\begin{attsatser}
    \item 
\end{attsatser}
\subsubsection*{Beslut}
\begin{attsatser}
    \item %förslaget till beslut bifalles
\end{attsatser}

\subsection{Val av protokolljusterare}

Protokolljusterare har till uppgift att kontrollera att protokollet i slutändan reflekterar de faktiska besluten och diskussionerna som fördes under mötet.
Utöver protokolljusteraren så ska mötesordförande och mötessekreteraren signera protokollet.
Vid styrelsemöten ska det endast vara en justerare.
Mötesordförande och mötessekreteraren kan inte vara justerare.

\subsubsection*{Förslag}
\begin{attsatser}
    \item 
\end{attsatser}
\subsubsection*{Beslut}
\begin{attsatser}
    \item %förslaget till beslut bifalles
\end{attsatser}

\section{Rapport} 
\subsection{Ordförande}
...

\subsection{Kassör}
...

\subsection{Vice-ordförande}
...

\subsection{SAMO}
...

\subsection{Mötessekreterare}
...

\newpage


\section{Diskussionspunkter}
\subsection*{Disk-punkt}
...

\subsection*{Disk-punkt}
...


\section{Besultspunkter}

\subsection{Beslut-punkt}
...
\begin{attsatser}
    \item 
\end{attsatser}
\subsubsection{Beslut}
\begin{attsatser}
    \item %Attsats 1 godkänns.
\end{attsatser}

\subsection{Beslut-punkt}
...
\begin{attsatser}
    \item 
\end{attsatser}
\subsubsection{Beslut} 
\begin{attsatser}     
\item %Attsats 1 bifalles. 
\end{attsatser}


\newpage
\section{Avslutande av möte}

\subsection{Nästa möte}
YYYY-MM-DD - HH:MM

\subsection{Mötets avslutande}
Mötet avslutades kl HH:MM 

\styrelsesignaturer

\end{document}
